% configurações
\documentclass[twocolumn]{article}

% português
\usepackage[utf8]{inputenc}
\usepackage[portuges,brazilian]{babel}
\usepackage[T1]{fontenc}

% pacotes
\usepackage{tikz}
\usepackage{amssymb,amsmath}
\usepackage[siunitx]{circuitikz}
\usepackage{hyperref}

% símbolos
\newcommand{\ohm}{\Omega}
\newcommand{\milli}{\mu}

\title{Entendendo Eletrônica}
\author{André Wagner}

\begin{document}

\maketitle

\section{Básico}

\subsection{Lei de Ohm}

\begin{description}
  \item[Voltagem ($V$)]: potencial. Equivalente à pressão d'água. Medido em Volts ($V$).
  \item[Corrente ($I$)]: quantidade de elétrons passando por um ponto em determinado momento. Equivalente ao fluxo de água. Medido em Amperes ($A$).
  \item[Resistência ($R$)]: o quanto o condutor se opõe à passagem de energia. Equivalente a um cano mais fino. Medido em Ohms ($\ohm$).
\end{description}

$$ V = IR $$
$$ I = \frac{V}{R} $$
$$ R = \frac{V}{I} $$

\url{http://www.falstad.com/circuit/e-ohms.html}

\begin{circuitikz} \draw
 (0,0) node[ground] {}
   to[short, l=$50\milli A$] (0,1)
   to[R, l=$100\ohm$] (0,3)
   -- (2,3)
   to[R, l=$1k\ohm$] (2,1)
   to[short, l=$5\milli A$] (2,0)
   node[ground] {}
 (1,3) to[short, -o] (1,4) node[anchor=south]{$+5V$}
;
\end{circuitikz}

No circuito acima, no lado esquerdo:
$$ V = IR \longrightarrow 5 = 100 \times 0,05 $$

No lado direito:
$$ V = IR \longrightarrow 5 = 1000 \times 0,005 $$

Isso significa que mais energia ``flui'' do lado esquerdo que do lado direito.


\subsection{Resistores}

\url{http://www.falstad.com/circuit/e-resistors.html}

Resistores se opõem à passagem de energia. Resistores em série tem a resistência somada.

\begin{circuitikz} \draw
 (0,0) to[R, l=$R_1$, o-] (2,0)
   to[R, l=$R_2$] (3.5,0)
   to[R, l=$R_n$, -o] (5.5,0)
;
\end{circuitikz}

$$ R = R_1 + R_2 + \cdots + R_n = \sum_{i=1}^{n} R_i $$

Resistores em paralelo tem a resistência distribuída.

\begin{circuitikz} \draw
 (0,1) to[short, o-] (1,1)
   to[short] (1,2)
   to[R, l=$R_1$] (3,2)
   to[short] (3,1)
   to[short, -o] (4,1)
 (1,1) to[R, l=$R_2$, *-*] (3,1)
 (1,1) to[short] (1,0)
   to[R, l=$R_n$] (3,0)
   to[short] (3,1)
;
\end{circuitikz}

$$ R = \frac{1}{R_1} + \frac{1}{R_2} + \cdots + \frac{1}{R_n} = \sum_{i=1}^{n} \frac{1}{R_i} $$

No caso especial de dois resistores em paralelo, a seguinte fórmula pode ser usada:

$$ \frac{R_1 R_2}{R_1 + R_2} $$

Exemplo:

\begin{circuitikz} \draw
 (0,0) to[R, l=$50\ohm (R_1)$, o-] (2,0)
   to[R, l=$150\ohm (R_2)$, *-*] (4,0)
 (2,0) to[short] (2,1)
   to[R, l=$200\ohm (R_3)$] (4,1)
   to[short] (4,0)
   to[short, -o] (5,0)
;
\end{circuitikz}

$$ R_1 + \Bigg(\frac{R_2 R_3}{R_2 + R_3}\Bigg) = 50 + \Bigg(\frac{150 \cdot 200}{150 + 200}\Bigg) = 135\ohm $$


\subsection{Capacitores}

Capacitores armazenam carga. À medida que a corrente flui para dentro do capacitor, a voltagem do capacitor aumenta. À medida que a voltagem se aproxima à voltagem da fonte, a corrente fluindo para dentro do capacidor diminui.

\url{http://www.falstad.com/circuit/e-cap.html}



\end{document}
